\chapter{Marco jurídico e institucional de la organización hospitalaria en España}
\label{ch:juridico}

% ═══════════════════════════════════════════════════════════
% SECCIÓN 2.1
% ═══════════════════════════════════════════════════════════

\section{Actualidad y persistencia del debate sobre la organización hospitalaria}
\label{sec:actualidad}

El debate sobre cómo debe organizarse la gestión hospitalaria pública en España no es un fenómeno reciente. Desde que la Ley 15/1997 habilitó de forma explícita la gestión de centros y servicios sanitarios públicos mediante cualesquiera entidades de naturaleza o titularidad pública admitidas en Derecho, el sistema hospitalario español ha acumulado casi tres décadas de experimentación institucional con fórmulas organizativas diversas: fundaciones sanitarias públicas, empresas públicas, consorcios, concesiones administrativas y distintas modalidades de concierto con el sector privado. Esta trayectoria no ha sido lineal: ha alternado fases de expansión de las formas de gestión indirecta con fases de reversión hacia la gestión directa, en un proceso marcado tanto por consideraciones de eficiencia como por cambios en las mayorías políticas autonómicas.

La actualidad de 2026 confirma que la organización hospitalaria sigue siendo un objeto de decisión política y jurídica de primer orden. Dos iniciativas normativas en tramitación ilustran este punto con claridad.

En primer lugar, el 10 de febrero de 2026, el Consejo de Ministros aprobó el anteproyecto de Ley de Gestión Pública e Integridad del Sistema Nacional de Salud \citep{ministerio2026gestion}. Según la exposición institucional del Ministerio de Sanidad, el anteproyecto tiene como objetivos principales garantizar la prioridad de la gestión pública directa en la provisión de servicios sanitarios, regular la gestión indirecta como modalidad excepcional sometida a requisitos reforzados de justificación y transparencia, y derogar expresamente la Ley 15/1997. Con independencia del resultado del proceso legislativo, la sola tramitación de una reforma de este alcance ---que pretende sustituir el marco habilitante que ha regido durante tres décadas--- evidencia que las reglas de organización hospitalaria no son un dato fijo del entorno institucional, sino una variable sometida a revisión periódica.

En segundo lugar, el Ministerio de Sanidad publicó en enero de 2026 el borrador del nuevo Estatuto Marco del personal estatutario de los servicios de salud \citep{ministerio2026estatuto}, con el propósito de actualizar el régimen establecido por la Ley 55/2003. El borrador ha generado un proceso de negociación con las organizaciones sindicales y profesionales médicas que, en la práctica, ha derivado en convocatorias de huelga y movilizaciones durante el primer trimestre de 2026. Desde la perspectiva de esta tesis, el interés de estos hechos reside en su contenido institucional más que en su dimensión mediática: el conflicto pone el foco en el régimen del personal sanitario, la dirección efectiva del trabajo clínico, la asignación de tiempos y guardias, y, en general, en cómo se estructura la autoridad dentro de las organizaciones hospitalarias. Estas son precisamente las variables que determinan la capacidad del gestor hospitalario para diseñar incentivos y modular el esfuerzo de los profesionales ---el núcleo del análisis económico de esta tesis.

A estas reformas en tramitación se suma la experiencia acumulada de las reversiones de concesiones administrativas en la Comunidad Valenciana. El modelo de concesión integral inaugurado con el Hospital de La Ribera (Alzira) en 1999 se extendió durante la década de 2000 a otros centros valencianos (Torrevieja, Dénia, Manises, Vinalopó). Sin embargo, a partir de 2018, la Generalitat Valenciana inició un proceso sistemático de reversión: el Hospital de La Ribera fue revertido en 2018, el Hospital de Torrevieja en 2021 y el Hospital de Dénia en 2024. El Hospital de Manises completó su reversión en 2024, mientras que la concesión del Hospital del Vinalopó fue prorrogada hasta 2030. Estas reversiones ---confirmadas judicialmente en varios casos, entre ellos por la STS 952/2021 relativa a Alzira--- constituyen evidencia empírica de que las formas de gestión hospitalaria no son irreversibles: son decisiones institucionales sujetas a revisión, lo que refuerza la pertinencia de evaluar sus consecuencias sobre el desempeño.

La conclusión de esta breve revisión es directa: si las reglas organizativas no importaran para el funcionamiento de los hospitales, no serían objeto recurrente de reforma legislativa ni de conflicto profesional. Esta tesis analiza precisamente en qué medida esas reglas ---la forma de gestión, el régimen del personal, el mecanismo de financiación--- afectan al desempeño hospitalario en sus distintas dimensiones.


% ═══════════════════════════════════════════════════════════
% SECCIÓN 2.2
% ═══════════════════════════════════════════════════════════

\section{Fundamento constitucional y marco legal básico}
\label{sec:fundamento}

\subsection{El derecho a la protección de la salud y el margen de configuración legislativa}

El artículo 43 de la Constitución Española de 1978 reconoce el derecho a la protección de la salud y establece que compete a los poderes públicos organizar y tutelar la salud pública a través de medidas preventivas y de las prestaciones y servicios necesarios. Este precepto se ubica entre los principios rectores de la política social y económica (Capítulo III del Título I), lo que implica que no configura un derecho fundamental directamente invocable ante los tribunales mediante recurso de amparo, sino un mandato al legislador y a los poderes públicos que, conforme al artículo 53.3 CE, informará la legislación positiva, la práctica judicial y la actuación de los poderes públicos.

La consecuencia analíticamente relevante es que el artículo 43 CE establece un fin ---la protección de la salud--- y atribuye a los poderes públicos la responsabilidad de organizar y tutelar, pero no constitucionaliza una única forma de organización de los servicios sanitarios. Existe un amplio margen de configuración legislativa: el legislador puede optar por modelos de gestión directa, indirecta o mixta, siempre que se garantice la efectividad del derecho y se respeten los principios de igualdad (art. 14 CE) y universalidad en el acceso. Esta apertura constitucional es la que ha permitido la diversidad organizativa que caracteriza al sistema hospitalario español y que esta tesis explota analíticamente.

\subsection{Distribución competencial y descentralización sanitaria}

La distribución de competencias en materia sanitaria entre el Estado y las comunidades autónomas se articula mediante dos preceptos constitucionales. El artículo 149.1.16.ª CE atribuye al Estado la competencia exclusiva sobre las bases y coordinación general de la sanidad, así como la legislación sobre productos farmacéuticos. El artículo 148.1.21.ª CE permite a las comunidades autónomas asumir competencias en sanidad e higiene, lo que se ha desarrollado en los respectivos Estatutos de Autonomía.

El resultado práctico ha sido un proceso de transferencias sanitarias que se inició con Cataluña en 1981 y se completó en enero de 2002 con la asunción de competencias por las diez comunidades autónomas que aún permanecían bajo la gestión directa del INSALUD (Instituto Nacional de la Salud). A partir de esa fecha, los diecisiete servicios autonómicos de salud ---más las ciudades autónomas de Ceuta y Melilla--- ejercen la función de planificación, financiación y provisión de servicios sanitarios en sus respectivos territorios \citep{observatorioHIT2024}. La disolución del INSALUD como gestor directo supuso el fin de un modelo centralizado de provisión y el inicio de una etapa de heterogeneidad institucional en la que cada comunidad autónoma ha podido configurar su red hospitalaria con arreglos organizativos propios, dentro del marco de la legislación básica estatal.

\subsection{Ley General de Sanidad (Ley 14/1986)}

La Ley 14/1986, de 25 de abril, General de Sanidad (BOE-A-1986-10499), constituye la norma fundacional del Sistema Nacional de Salud. Sus preceptos más relevantes para el análisis de la organización hospitalaria son los siguientes.

El artículo 44 define el Sistema Nacional de Salud como el conjunto de los servicios de salud de la Administración del Estado y de las comunidades autónomas, integrados y coordinados bajo los principios de universalidad, equidad, calidad y participación. El artículo 56 define el hospital como el establecimiento encargado tanto del internamiento clínico como de la asistencia especializada y complementaria que requiera su zona de influencia. Los artículos 65 a 67 estructuran la organización territorial en áreas de salud, cada una con al menos un hospital general, y prevén la vinculación de hospitales privados al sistema público mediante distintas fórmulas. El artículo 90 regula la contratación de servicios con hospitales privados mediante conciertos, estableciendo que los centros concertados deben cumplir los mismos requisitos de calidad, control y accesibilidad que los centros públicos.

La Ley General de Sanidad configuró así un sistema de provisión predominantemente pública, pero abierto desde su origen a la participación del sector privado mediante conciertos, un rasgo que ha condicionado toda la evolución posterior \citep{peman2005}.

\subsection{Ley de Cohesión y Calidad (Ley 16/2003)}

La Ley 16/2003, de 28 de mayo, de Cohesión y Calidad del Sistema Nacional de Salud (BOE-A-2003-10715), fue la respuesta legislativa a la necesidad de articular un marco de coordinación tras la culminación de las transferencias sanitarias. Sus principales aportaciones son:

El artículo 7 establece un catálogo de prestaciones del SNS que incluye la salud pública, la atención primaria, la atención especializada, la atención sociosanitaria, la atención de urgencias y las prestaciones farmacéuticas. El artículo 8 regula la cartera común de servicios, distinguiendo entre básica, suplementaria y de servicios accesorios. Esta definición amplia de las prestaciones ---que va más allá del internamiento--- tiene una consecuencia analítica directa: el producto hospitalario no se agota en el número de altas o estancias, sino que incluye dimensiones de calidad, adecuación diagnóstica y continuidad asistencial que son difícilmente reducibles a un indicador único.

Los artículos 69 a 73 regulan el Consejo Interterritorial del Sistema Nacional de Salud, órgano de coordinación integrado por la persona titular del Ministerio de Sanidad y los consejeros autonómicos responsables de sanidad, con funciones de coordinación, cooperación y definición de posiciones comunes. El artículo 63 crea el Observatorio del Sistema Nacional de Salud, con la función de proporcionar un análisis permanente del SNS mediante estudios comparados de los servicios de salud autonómicos en ámbitos que incluyen la organización, la provisión, la gestión y los resultados. Esta previsión legal de evaluación comparada refuerza la legitimidad de una tesis que se propone evaluar el desempeño hospitalario con enfoque empírico.

\subsection{Ley 15/1997: habilitación de nuevas formas de gestión}

La Ley 15/1997, de 25 de abril, sobre habilitación de nuevas formas de gestión del Sistema Nacional de Salud (BOE-A-1997-9021), es una norma breve ---un artículo único y una disposición adicional--- pero de extraordinaria importancia para el objeto de esta tesis. Su artículo único establece:

\begin{quote}
\textit{La gestión y administración de los centros, servicios y establecimientos sanitarios de protección de la salud o de atención sanitaria o sociosanitaria podrá llevarse a cabo directamente o indirectamente a través de cualesquiera entidades de naturaleza o titularidad pública admitidas en Derecho.}
\end{quote}

La formulación es deliberadamente abierta: no enumera un catálogo cerrado de formas jurídicas, sino que habilita cualquier entidad de naturaleza o titularidad pública. Esta apertura proporcionó la base legal para la proliferación de formas de gestión que se produjo en los años siguientes: fundaciones públicas sanitarias, empresas públicas, consorcios, sociedades mercantiles de capital público y, mediante interpretación extensiva, concesiones administrativas a entidades privadas para la gestión de centros sanitarios públicos.

\subsection{Real Decreto 29/2000: desarrollo reglamentario}

El Real Decreto 29/2000, de 14 de enero, sobre nuevas formas de gestión del INSALUD, desarrolló reglamentariamente la habilitación de la Ley 15/1997 en el ámbito del entonces existente INSALUD. El RD articuló cuatro figuras organizativas específicas: fundaciones públicas sanitarias, consorcios, sociedades estatales y otras entidades de Derecho público. Para cada una estableció reglas sobre órganos de gobierno, régimen patrimonial, contratación de personal y mecanismos de control. Aunque el RD 29/2000 fue concebido para el INSALUD (disuelto en 2002 tras las transferencias), sus categorías organizativas sirvieron de referencia para las comunidades autónomas en el diseño de sus propios modelos de gestión.

\subsection{Ley 55/2003: Estatuto Marco del personal estatutario}

La Ley 55/2003, de 16 de diciembre, del Estatuto Marco del personal estatutario de los servicios de salud (BOE-A-2003-23101), regula la relación de empleo del personal estatutario como una relación funcionarial de carácter especial. La Ley establece las normas sobre acceso (mediante oferta pública de empleo y procedimientos de selección), provisión de plazas, retribuciones (básicas y complementarias), jornada, permisos, situaciones administrativas y régimen disciplinario. La importancia de esta norma para el análisis de la tesis se desarrolla en la Sección~\ref{sec:personal}.

\subsection{Síntesis}

El ordenamiento configura un sistema sanitario con fines públicos ---protección de la salud, universalidad, equidad--- pero formas jurídicas plurales para la organización de la provisión. La Constitución no constitucionaliza un único modelo de gestión; la legislación básica habilita la diversidad organizativa; y las comunidades autónomas disponen de amplias competencias para diseñar sus redes hospitalarias. La STC 84/2015, al resolver un recurso de inconstitucionalidad contra disposiciones de la legislación sanitaria madrileña, confirmó que la opción por la gestión indirecta de servicios sanitarios no es per se inconstitucional, siempre que se respeten las bases estatales y los principios de igualdad y universalidad del sistema.

Esta pluralidad organizativa ---habilitada normativamente, no accidental--- es la que genera la variación institucional que la tesis explota analíticamente. El siguiente paso es describir las formas concretas que esa pluralidad ha adoptado en la práctica.


% ═══════════════════════════════════════════════════════════
% SECCIÓN 2.3
% ═══════════════════════════════════════════════════════════

\section{Tipología de las formas de gestión y contratación hospitalaria}
\label{sec:tipologia}

\subsection{Gestión directa}
\label{sec:gestion_directa}

\subsubsection{El hospital integrado en el servicio de salud autonómico}

La forma organizativa predominante en el sistema hospitalario español es el hospital público integrado directamente en el servicio de salud de la comunidad autónoma correspondiente. Estos centros carecen de personalidad jurídica propia: son unidades orgánicas del servicio autonómico de salud, sujetas íntegramente al Derecho público en materia de contratación, presupuestación, gestión de personal y control del gasto. Su personal es estatutario (Ley 55/2003), seleccionado mediante oferta pública de empleo, con retribuciones fijadas en mesas sectoriales autonómicas y sujeto a intervención previa del gasto por la intervención general de la comunidad autónoma.

Desde la perspectiva del análisis de incentivos, el hospital integrado representa el extremo del gradiente de autonomía: la dirección del centro dispone de un margen muy limitado para modular las condiciones de trabajo o la retribución del personal médico. Las decisiones sobre plantillas, retribuciones complementarias, carrera profesional y selección de personal se adoptan a un nivel superior al de la dirección hospitalaria.

\subsubsection{Entes instrumentales con personalidad jurídica propia}

Dentro de la gestión directa ---entendida como gestión mediante entidades de titularidad pública--- existen formas organizativas con personalidad jurídica propia y distintos grados de autonomía. Las principales son:

\textit{Organismos autónomos y entidades públicas empresariales.} El caso más relevante es Osakidetza (Servicio Vasco de Salud), configurada como ente público de Derecho privado, lo que le permite una mayor flexibilidad en la contratación y gestión del personal, aunque la práctica ha convergido hacia un régimen próximo al estatutario. En Cataluña, el Institut Català de la Salut (ICS) es una empresa pública adscrita al Departament de Salut que gestiona directamente los hospitales y centros de atención primaria de mayor tamaño.

\textit{Consorcios sanitarios.} Cataluña ha desarrollado un modelo organizativo singular basado en consorcios sanitarios ---entidades de Derecho público con participación de diversas administraciones y, en algunos casos, entidades privadas sin ánimo de lucro--- que gestionan hospitales y otros centros asistenciales. La Ley 15/1990, de Ordenación Sanitaria de Cataluña, configuró un sistema plural de provisión articulado en torno al Servei Català de la Salut (CatSalut) como entidad aseguradora-compradora, que contrata servicios con una red diversa de proveedores: el ICS (gestión directa), consorcios, fundaciones, hospitales municipales, entidades religiosas y centros privados concertados. El sistema se articula a través del SISCAT (Sistema Sanitari Integral d'Utilització Pública de Catalunya), que integra a todos los proveedores acreditados.

\textit{Fundaciones públicas sanitarias.} El Real Decreto 29/2000 articuló las fundaciones públicas sanitarias como una de las nuevas formas de gestión del INSALUD. En Galicia, cinco hospitales fueron constituidos como fundaciones públicas sanitarias durante los años 2000, con personal laboral y mayor autonomía de gestión. Sin embargo, entre 2007 y 2008, la Xunta de Galicia procedió a la extinción de estas fundaciones y a la integración de sus centros en la red del SERGAS (Servizo Galego de Saúde), con la consiguiente estatutarización del personal.

\textit{Empresas públicas.} En Andalucía, varias Agencias Públicas Empresariales Sanitarias (APES) gestionaron hospitales y otros centros durante más de una década con personal laboral sujeto a convenio colectivo propio. A partir de 2019, la Junta de Andalucía inició un proceso de liquidación de estas agencias y de integración de sus centros y personal en el Servicio Andaluz de Salud (SAS), proceso que se completó entre 2020 y 2022.

Estos ejemplos autonómicos ilustran un patrón relevante para el análisis: las formas de gestión con mayor autonomía (fundaciones, empresas públicas) han experimentado en varios casos un proceso de reversión hacia la gestión integrada con personal estatutario, lo que indica que la tendencia institucional predominante ha sido hacia la homogeneización del régimen de gestión directa.

\begin{table}[htbp]
\centering
\caption{Formas de gestión directa en el sistema hospitalario español}
\label{tab:gestion_directa}
\small
\begin{tabular}{p{3.2cm}p{3cm}p{2.8cm}p{3.5cm}}
\toprule
\textbf{Naturaleza jurídica} & \textbf{Régimen de personal} & \textbf{Grado de autonomía} & \textbf{Ejemplos} \\
\midrule
Hospital integrado en servicio de salud & Estatutario & Muy limitada & Mayoría de hospitales SNS \\[4pt]
Organismo autónomo / EPE & Estatutario o laboral & Media & Osakidetza, ICS \\[4pt]
Consorcio sanitario & Laboral (convenio propio) & Media-alta & Consorcios catalanes \\[4pt]
Fundación pública sanitaria & Laboral & Alta & Fundaciones gallegas (extintas) \\[4pt]
Empresa pública & Laboral (convenio propio) & Alta & APES andaluzas (liquidadas) \\[4pt]
Sociedad mercantil pública & Laboral & Alta & Casos puntuales \\
\bottomrule
\end{tabular}
\end{table}

\subsection{Gestión indirecta y contratación con el sector privado}
\label{sec:gestion_indirecta}

\subsubsection{Concesión administrativa}

La concesión administrativa de servicios sanitarios ha adoptado en España dos modalidades principales que es imprescindible distinguir analíticamente.

La primera es la \textit{iniciativa de financiación privada} (PFI, por sus siglas en inglés), en la que una empresa privada construye y financia la infraestructura hospitalaria y gestiona los servicios no clínicos (mantenimiento, limpieza, restauración, logística), mientras que la asistencia sanitaria es prestada por personal estatutario del servicio de salud. La Comunidad de Madrid aplicó este modelo en la construcción de siete nuevos hospitales entre 2005 y 2008. En este esquema, la dirección clínica sigue siendo pública y el régimen del personal es estatutario, por lo que la autonomía del gestor para diseñar incentivos del personal médico no difiere sustancialmente de la del hospital integrado.

La segunda modalidad es la \textit{concesión integral}, cuyo modelo paradigmático es el Hospital de La Ribera (Alzira, Valencia), inaugurado en 1999. En la concesión integral, el concesionario privado asume la gestión completa del hospital ---incluida la asistencia clínica--- a cambio de un pago capitativo por habitante asignado al área de salud. El concesionario selecciona y contrata al personal con régimen laboral, diseña autónomamente los esquemas retributivos internos y asume el riesgo económico de la gestión. La Comunidad Valenciana extendió este modelo a otros cuatro hospitales: Torrevieja (2006), Dénia (2009), Manises (2009) y Vinalopó (2010).

Como se ha señalado en la Sección~\ref{sec:actualidad}, la Comunidad Valenciana inició a partir de 2018 un proceso de reversión de estas concesiones. El Hospital de La Ribera fue revertido en abril de 2018 al expirar la concesión; el Hospital de Torrevieja fue revertido en octubre de 2021, tras un proceso contencioso confirmado judicialmente (STS 952/2021); el Hospital de Dénia fue revertido en 2024; y el Hospital de Manises completó su reversión en 2024. La concesión del Hospital del Vinalopó fue prorrogada hasta 2030. En la Comunidad de Madrid, además de los hospitales PFI, se desarrollaron cuatro concesiones con componente de gestión clínica, aunque con modelos heterogéneos.

Desde la perspectiva del modelo de agencia de esta tesis, la concesión integral es el caso empíricamente más limpio de estructura descentralizada: el financiador público contrata con el concesionario (hospital) para un nivel de servicio, y el concesionario diseña autónomamente los contratos del personal médico. El pago capitativo, además, genera incentivos específicos sobre la eficiencia en intensidad (contención del coste por episodio) que difieren de los incentivos generados por el pago por actividad.

\subsubsection{Concierto sanitario}

El concierto sanitario es el instrumento clásico de contratación de servicios hospitalarios con el sector privado, previsto en los artículos 67 y 90 de la Ley General de Sanidad. Mediante el concierto, el servicio de salud autonómico contrata con un hospital privado la prestación de servicios asistenciales a pacientes del SNS, a cambio de un precio pactado por módulos de actividad (precio por estancia, por alta, por procedimiento o por consulta).

El hospital concertado mantiene plena autonomía sobre la contratación y retribución de su personal ---que se rige por el Derecho laboral ordinario--- y gestiona internamente sus recursos de acuerdo con su propia estructura organizativa. Sin embargo, su actividad para el SNS está sujeta a los módulos de precio fijados por la administración sanitaria y a los controles de calidad y accesibilidad que establezca el concierto. La red de hospitales concertados es particularmente extensa en Cataluña (donde forma parte del SISCAT) y en la Comunidad de Madrid.

\subsubsection{Acción concertada}

Algunas comunidades autónomas han desarrollado en los últimos años una modalidad de contratación denominada \textit{acción concertada}, que se diferencia del concierto clásico en aspectos formales y procedimentales. La acción concertada se concibe como un régimen específico para la contratación de servicios sanitarios con entidades privadas, con requisitos propios de acreditación, precios regulados y obligaciones de información. Aunque su extensión es aún limitada, representa una evolución del marco contractual entre el sector público y el privado en materia sanitaria.

\subsection{Síntesis analítica}
\label{sec:sintesis_tipologia}

La descripción precedente no pretende agotar la taxonomía de formas de gestión hospitalaria en España, sino identificar la variable institucional más relevante para el análisis económico de esta tesis: el \textit{gradiente de autonomía} del gestor hospitalario para diseñar la estructura interna de incentivos.

En un extremo del gradiente se sitúa el hospital público integrado, con personal estatutario, retribuciones centralizadas y escasa autonomía del gestor para modular incentivos individuales. En el otro extremo se sitúan las concesiones integrales y los hospitales privados de mercado, donde la dirección del centro selecciona, contrata y retribuye al personal con plena autonomía, sujeta únicamente al marco laboral general. Entre ambos extremos se ubican las fundaciones, consorcios, empresas públicas y hospitales concertados, con grados intermedios de autonomía que dependen de su naturaleza jurídica, del convenio colectivo aplicable y de la práctica de gestión.

Las variables institucionales analíticamente relevantes para el modelo de agencia de esta tesis son cuatro:

\begin{enumerate}[label=(\roman*)]
  \item \textit{Autonomía del gestor para diseñar contratos internos}: determina si el hospital puede vincular la retribución del médico al desempeño y en qué medida.
  \item \textit{Régimen de personal} (estatutario versus laboral): condiciona los instrumentos de contratación, selección y retribución disponibles.
  \item \textit{Reparto de riesgos entre financiador y proveedor}: un hospital financiado por presupuesto global no soporta riesgo de actividad; un hospital con pago capitativo soporta riesgo económico completo.
  \item \textit{Capacidad de supervisión del financiador}: los mecanismos de control varían desde la intervención previa del gasto (hospitales integrados) hasta el control ex post mediante indicadores contractuales (concesiones y conciertos).
\end{enumerate}

La variable empírica $D^{desc}$ de la estimación econométrica (Capítulo~\ref{ch:datos}) operacionaliza el gradiente de autonomía mediante una clasificación dicotómica que distingue hospitales centralizados (gestión directa con personal estatutario) de hospitales descentralizados (gestión indirecta o privada con mayor autonomía contractual). El Diseño B refina esta clasificación en tres categorías con anclaje económico más preciso.


% ═══════════════════════════════════════════════════════════
% SECCIÓN 2.4
% ═══════════════════════════════════════════════════════════

\section{Régimen del personal, autonomía de gestión y gobernanza interna}
\label{sec:personal}

\subsection{El personal estatutario: marco general}

La Ley 55/2003, del Estatuto Marco del personal estatutario de los servicios de salud (BOE-A-2003-23101), regula la relación de empleo del personal que presta servicios en los centros e instituciones sanitarias de los servicios de salud de las comunidades autónomas. Esta relación se configura como una relación funcionarial de carácter especial, distinta tanto de la relación funcionarial general (regulada por el EBEP) como de la relación laboral ordinaria (regulada por el Estatuto de los Trabajadores).

El acceso al empleo estatutario se realiza mediante oferta pública de empleo (OPE) y procedimientos de selección (oposición, concurso-oposición o concurso) convocados por las comunidades autónomas, con plazos que en la práctica pueden extenderse durante varios años desde la convocatoria hasta la toma de posesión. Esta rigidez temporal en la selección limita la capacidad de la dirección hospitalaria para adaptar las plantillas a las necesidades asistenciales cambiantes.

\subsection{Estructura retributiva y margen para incentivos}

La estructura retributiva del personal estatutario se compone de retribuciones básicas ---sueldo base y trienios, fijados a nivel estatal--- y retribuciones complementarias, entre las que se incluyen el complemento de destino, el complemento específico y el complemento de productividad variable. Este último es, en principio, el instrumento disponible para vincular parte de la retribución al desempeño individual o colectivo.

Sin embargo, en la práctica, el margen para incentivos es estrecho. El complemento de productividad variable representa, según la comunidad autónoma, entre el 1,5\% y el 15\% de la retribución total del médico especialista. La fijación de los criterios de productividad se negocia en las mesas sectoriales autonómicas, no a nivel de hospital, y los indicadores utilizados suelen ser agregados (cumplimiento de objetivos del contrato-programa del centro) más que individuales. El resultado es un esquema retributivo con potencia de incentivos muy limitada: la mayor parte de la retribución del médico estatutario es fija e independiente de su desempeño individual \citep{monereo2007}.

La carrera profesional, regulada en el artículo 40 de la Ley 55/2003 y desarrollada de forma heterogénea por las comunidades autónomas, constituye el otro instrumento de incentivos a largo plazo. La carrera profesional reconoce el desarrollo profesional mediante niveles o grados asociados a complementos retributivos, y exige la evaluación de competencias, formación y actividad asistencial. Sin embargo, su implantación ha sido desigual y, en varias comunidades, los avances han sido limitados o se han retrasado significativamente.

\subsection{Personal laboral en entes instrumentales}

En los hospitales gestionados por fundaciones públicas, empresas públicas, consorcios u otras entidades con personalidad jurídica propia, el personal se ha regido históricamente por el Derecho laboral, mediante contratos laborales sujetos a convenio colectivo propio de la entidad. Este régimen permite a la dirección del hospital una mayor flexibilidad en la contratación (selección directa, no sujeta a OPE), en la retribución (posibilidad de vincular una parte significativa del salario al desempeño individual o por unidades) y en la organización del trabajo (mayor capacidad para reasignar funciones y modificar condiciones).

No obstante, la tendencia institucional predominante en la última década ha sido la estatutarización del personal de estos entes instrumentales, como ilustran los casos de las fundaciones gallegas (integradas en el SERGAS entre 2007 y 2008, con posterior estatutarización del personal) y de las APES andaluzas (liquidadas entre 2019 y 2022, con integración del personal en el régimen estatutario del SAS). Esta convergencia hacia el régimen estatutario reduce progresivamente el universo de hospitales públicos con autonomía retributiva plena.

\subsection{Autonomía del gestor hospitalario}

La autonomía del gestor hospitalario ---entendida como la capacidad efectiva del director gerente para tomar decisiones sobre la organización interna de recursos humanos, la retribución del personal y la asignación de medios--- varía sustancialmente según la forma de gestión.

En los hospitales públicos integrados, la autonomía del gestor es muy limitada. Las plantillas se fijan centralizadamente por el servicio de salud autonómico; las retribuciones se negocian en mesas sectoriales a nivel autonómico; la selección de personal se realiza mediante OPE con plazos que escapan al control de la dirección del centro; y la contratación temporal está sujeta a restricciones presupuestarias y administrativas. La dirección del hospital puede modular en cierta medida la organización funcional interna (distribución de guardias, asignación de tareas, programación quirúrgica), pero no puede alterar significativamente la estructura retributiva del personal médico.

En las concesiones integrales y en los hospitales privados (concertados o de mercado), la dirección dispone de plena autonomía para seleccionar personal, fijar retribuciones, diseñar esquemas de incentivos variables y reorganizar la estructura funcional del centro. Esta autonomía está limitada, en el caso de los hospitales concertados, por los módulos de precio del concierto y por los indicadores de calidad y accesibilidad exigidos por la administración sanitaria.

\subsection{Directivos de hospital}

El régimen de los directivos de los hospitales del SNS es heterogéneo y presenta debilidades reconocidas. La designación del director gerente es, en la mayoría de las comunidades autónomas, un nombramiento de libre designación por la autoridad sanitaria, lo que introduce un componente de discrecionalidad política en la selección. La Ley 55/2003 (art. 26) establece que el personal directivo será nombrado atendiendo a los principios de igualdad, mérito, capacidad y publicidad, pero la práctica efectiva dista con frecuencia de estos principios. La rotación frecuente de los directivos hospitalarios ---asociada a los ciclos políticos autonómicos--- limita la capacidad de implementar reformas organizativas de medio plazo.

\subsection{Conexión con el modelo teórico}

La descripción del régimen de personal tiene una traducción directa en el modelo de agencia de la tesis. En un hospital con personal estatutario, la dirección no puede modular significativamente los incentivos del médico: las retribuciones están fijadas centralizadamente, el complemento variable es marginal y la selección de personal escapa al control del gestor. Esta configuración se aproxima a la \textit{estructura centralizada} del modelo teórico (Capítulo~\ref{ch:modelo}): el principal (financiador) contrata de facto tanto con el hospital como con el médico, porque las condiciones esenciales del empleo médico se determinan a un nivel superior al de la dirección hospitalaria.

En un hospital con personal laboral o gestión concesional, la dirección diseña autónomamente los contratos del personal médico: selecciona, retribuye y puede vincular una parte sustancial de la retribución al desempeño individual o por unidades funcionales. Esta configuración se aproxima a la \textit{estructura descentralizada}: el principal contrata solo con el hospital, y es el hospital quien diseña el contrato del médico.

El borrador del nuevo Estatuto Marco publicado en enero de 2026 \citep{ministerio2026estatuto} pone de manifiesto que este régimen está en proceso de revisión, lo que refuerza la relevancia del análisis: las reglas del personal no son un dato fijo, sino una variable institucional que puede alterarse mediante reformas normativas.


% ═══════════════════════════════════════════════════════════
% SECCIÓN 2.5
% ═══════════════════════════════════════════════════════════

\section{Financiación hospitalaria y verificabilidad del output}
\label{sec:financiacion}

\subsection{Presupuesto global retrospectivo}

El mecanismo de financiación históricamente dominante en los hospitales públicos del SNS ha sido el presupuesto global retrospectivo: el hospital recibe una asignación presupuestaria anual basada fundamentalmente en el gasto histórico, con ajustes incrementales que reflejan la inflación, las decisiones de inversión y las modificaciones de plantilla. Este mecanismo equivale, en la terminología de la teoría de incentivos, a un pago fijo desvinculado del output: el volumen de financiación que recibe el hospital es esencialmente independiente de su actividad asistencial o de la calidad de sus resultados.

La debilidad de este esquema como instrumento de incentivos fue identificada tempranamente. El denominado Informe Abril (1991), elaborado por una comisión de análisis y evaluación del SNS, señaló entre sus principales conclusiones la necesidad de vincular la financiación hospitalaria a la actividad y los resultados, superando el modelo de presupuesto histórico. Esta recomendación fue el antecedente inmediato de la introducción de los contratos-programa.

\subsection{Contratos-programa}

El contrato-programa es el principal instrumento de relación entre el servicio de salud autonómico y el hospital público. Introducidos a partir de 1993 por el INSALUD y adoptados posteriormente por los servicios autonómicos de salud, los contratos-programa establecen un marco de objetivos y compromisos entre el financiador y el centro hospitalario.

Un contrato-programa típico especifica: objetivos de actividad (número de altas, estancias, consultas, intervenciones), objetivos de calidad (indicadores de seguridad clínica, tiempos de espera, satisfacción del paciente), objetivos de eficiencia (estancia media, coste por alta, índice de rotación) y, en algunos casos, objetivos de docencia e investigación. El grado de cumplimiento de estos objetivos se vincula al sistema de Dirección por Objetivos (DPO) del centro, que a su vez determina la cuantía del complemento de productividad variable del personal.

El contrato-programa introduce, por tanto, un elemento de pago variable en un esquema que sigue siendo predominantemente de pago fijo. Sin embargo, la potencia efectiva de los incentivos depende de la proporción de la financiación que esté genuinamente vinculada al cumplimiento de objetivos y de la credibilidad de las consecuencias del incumplimiento. En la práctica, el presupuesto base del hospital rara vez se reduce significativamente por incumplimiento de objetivos, lo que limita la potencia del mecanismo.

\subsection{Pago por actividad: GRD y sistemas de clasificación de pacientes}

Los Grupos Relacionados por el Diagnóstico (GRD) son un sistema de clasificación de pacientes que agrupa los episodios de hospitalización en categorías clínicamente coherentes y de consumo de recursos homogéneo. El sistema permite comparar la actividad hospitalaria ajustada por complejidad (case-mix) y, potencialmente, vincular la financiación a la actividad ponderada.

En España, el Conjunto Mínimo Básico de Datos (CMBD) fue implantado en 1987 como registro normalizado de los episodios de hospitalización, incluyendo los diagnósticos y procedimientos codificados que permiten la asignación de GRD. El Real Decreto 69/2015 amplió este registro al Registro de Actividad de Atención Especializada (RAE-CMBD), que incluye información sobre hospitalización, cirugía ambulatoria, hospital de día y procedimientos ambulatorios especiales.

En la mayoría de las comunidades autónomas, los GRD se utilizan principalmente como instrumento de gestión interna y de benchmarking entre hospitales, no como mecanismo de pago directo. Cataluña es la comunidad autónoma donde el pago por actividad basado en GRD está más desarrollado: CatSalut utiliza los GRD como base para el cálculo de los pagos a los hospitales de la red SISCAT, tanto públicos como concertados. En el resto de comunidades, el peso de los GRD en la financiación efectiva es menor, aunque su uso en la evaluación del desempeño hospitalario es generalizado.

\subsection{Módulos de concierto}

Los hospitales privados que prestan servicios al SNS mediante concierto sanitario se financian a través de módulos de precio fijados por las comunidades autónomas. Estos módulos establecen tarifas por unidad de actividad ---precio por estancia, por alta, por procedimiento quirúrgico, por consulta externa--- que constituyen la base económica del concierto. Los módulos de concierto representan un mecanismo de pago por actividad que genera incentivos directos sobre el volumen de producción: el hospital recibe más financiación cuanta más actividad realice, dentro de los límites que establezca el concierto.

\subsection{Pago capitativo}

El pago capitativo consiste en una tarifa fija por habitante asignado al área de salud del hospital, independientemente de la actividad efectivamente realizada. Este mecanismo se ha utilizado en las concesiones administrativas de tipo Alzira, donde el concesionario recibe un pago anual por habitante ajustado que cubre la totalidad de la asistencia sanitaria (hospitalaria y, en algunos modelos, también primaria) de la población asignada.

El pago capitativo genera incentivos radicalmente distintos al pago por actividad: el hospital tiene incentivos para contener el coste por episodio (eficiencia en intensidad) y para reducir la utilización innecesaria, puesto que toda actividad adicional no genera ingreso marginal sino coste adicional. Este mecanismo también opera en el ámbito del mutualismo administrativo (MUFACE, MUGEJU, ISFAS), donde las aseguradoras privadas que cubren a los funcionarios reciben un pago capitativo por persona asegurada.

\subsection{Conexión con el modelo teórico y la variable empírica}

El problema analítico central que emerge de la descripción de los mecanismos de financiación es que las señales de output disponibles para el financiador no son igualmente precisas para todas las dimensiones del producto hospitalario.

El volumen de actividad ---número de altas, estancias, intervenciones, consultas--- es relativamente fácil de medir, registrar y verificar. Los sistemas de información hospitalaria (SIAE, CMBD/RAE-CMBD) proporcionan datos exhaustivos y normalizados sobre la actividad asistencial. La intensidad de la atención ---la adecuación de los procedimientos diagnósticos y terapéuticos a la situación clínica de cada paciente, la profundidad del estudio diagnóstico, la pertinencia de las indicaciones terapéuticas--- es mucho más costosa de observar, más difícil de codificar en indicadores contractuales verificables y más susceptible de manipulación mediante upcoding o selección de pacientes.

En la notación del modelo teórico (Capítulo~\ref{ch:modelo}), esta asimetría en la verificabilidad se captura mediante las varianzas de las señales de output: la señal de cantidad ($\sigma_1^2$) es relativamente precisa, mientras que la señal de intensidad ($\sigma_2^2$) es ruidosa. La consecuencia, bien establecida en la teoría de multitarea \citep{holmstrom1991}, es que vincular los incentivos a la señal más precisa puede generar un desplazamiento de esfuerzo desde la dimensión menos verificable.

La variable empírica $pct^{SNS}$ de la estimación (Capítulo~\ref{ch:datos}) captura parcialmente esta dimensión: un hospital con alta financiación pública ($pct^{SNS}$ alto) opera predominantemente bajo instrumentos que enfatizan el volumen de actividad (contratos-programa, GRD) y que miden con mayor precisión la cantidad que la intensidad. Un hospital con baja financiación pública opera bajo esquemas de pago directo o de seguros privados que pueden incorporar señales diferentes sobre la calidad y la intensidad de la atención.


% ═══════════════════════════════════════════════════════════
% SECCIÓN 2.6 — TEXTO LITERAL (NO MODIFICAR)
% ═══════════════════════════════════════════════════════════

\section{Del lenguaje jurídico de la organización sanitaria al lenguaje económico de la agencia}
\label{sec:bisagra}

Los capítulos anteriores de esta tesis han descrito, con anclaje normativo, las formas de gestión habilitadas por el ordenamiento, las diferencias de control entre ellas, el régimen del personal y las fórmulas de financiación y contratación. El propósito de esta sección es explicitar cómo esas categorías jurídico-institucionales se traducen en los parámetros de un problema de análisis económico bajo información asimétrica.

\subsection{El hospital como organización compleja con múltiples agentes}

El Derecho sanitario organiza la provisión hospitalaria mediante una cadena de delegaciones sucesivas. La administración sanitaria autonómica (financiador) fija el marco presupuestario, la cartera de servicios y los instrumentos de control. La dirección del hospital recibe un mandato de gestión ---más o menos amplio según la forma jurídica del centro--- y organiza internamente los recursos humanos, materiales y tecnológicos. Los profesionales clínicos ---fundamentalmente los médicos--- toman decisiones asistenciales con alto contenido técnico y elevada discrecionalidad, en un contexto donde la información relevante (estado del paciente, idoneidad de los procedimientos, esfuerzo efectivamente dedicado) es en gran medida local y no transmisible íntegramente a los niveles superiores de la cadena.

La teoría económica de la agencia \citep{laffont2002} ofrece un lenguaje parsimonioso para formalizar esta estructura. En la terminología de los modelos principal--agente, el financiador público actúa como \textit{principal}: define objetivos (salud poblacional, accesibilidad, eficiencia) y diseña instrumentos (presupuestos, contratos-programa, módulos de concierto, indicadores) para inducir al agente a comportarse de acuerdo con esos objetivos. El hospital actúa como primer \textit{agente}: recibe recursos y un mandato, pero dispone de información privada sobre sus costes reales, su capacidad organizativa y su esfuerzo efectivo. El médico actúa como segundo agente ---o subagente en una estructura de delegación secuencial---: su esfuerzo determina directamente la cantidad y la calidad de la atención prestada, pero ese esfuerzo no es plenamente observable ni por la dirección del hospital ni, menos aún, por el financiador.

Lo que distingue esta cadena de un problema de agencia estándar es la \textit{multiplicidad de outputs}. El hospital no produce un bien homogéneo. Produce, al menos, un volumen de atención (número de pacientes atendidos, altas, intervenciones) y una intensidad de atención por episodio (recursos diagnósticos y terapéuticos aplicados a cada paciente). Estas dos dimensiones no son independientes en la función de producción: el tiempo y el esfuerzo que el médico dedica a ver más pacientes pueden competir con el tiempo dedicado a aplicar una intensidad diagnóstica adecuada a cada uno. Esta interacción entre tareas ---que puede ser de sustitución o de complementariedad según el tipo de servicio--- es el mecanismo central del modelo teórico que se desarrolla en el Capítulo~\ref{ch:modelo}.

\subsection{Estructura organizativa como estructura de agencia}

Las formas de gestión descritas en la Sección~\ref{sec:tipologia} no son, desde esta perspectiva, meras categorías administrativas. Cada forma de gestión configura una \textit{estructura de agencia} distinta, que determina quién diseña el contrato del médico, qué instrumentos están disponibles para alinear incentivos y qué restricciones operan sobre la capacidad de supervisión y retribución.

En un hospital público de gestión directa con personal estatutario, el financiador y la administración sanitaria intervienen simultáneamente en la fijación de las condiciones del personal médico: la selección se realiza por oposición centralizada, las retribuciones se determinan en mesas sectoriales autonómicas, y la capacidad de la dirección del hospital para modular incentivos individuales se limita al estrecho margen del complemento de productividad variable. En términos del modelo de agencia, esta configuración se aproxima a una \textit{estructura centralizada}: el principal (financiador) contrata, de facto, tanto con el hospital como con el médico, porque las condiciones esenciales del empleo médico se fijan a un nivel superior al de la dirección hospitalaria.

En un hospital gestionado mediante concesión administrativa, fundación con personal laboral o empresa privada concertada, la situación es cualitativamente distinta. El financiador contrata con la entidad gestora para un determinado nivel de actividad y, en su caso, de calidad. La entidad gestora diseña autónomamente los contratos de sus profesionales: selecciona, retribuye y puede vincular una parte significativa de la retribución al desempeño individual o por unidades. En términos del modelo de agencia, esta configuración se aproxima a una \textit{estructura descentralizada}: el principal contrata solo con el hospital, y es el hospital quien diseña el contrato del médico.

Esta correspondencia no es perfecta ---ninguna traducción entre categorías jurídicas y categorías económicas lo es---, pero captura la variable institucional más relevante para el análisis: el grado de autonomía efectiva del gestor hospitalario para diseñar la estructura interna de incentivos. La variable empírica $D^{desc}$ de la estimación econométrica (Capítulo~\ref{ch:datos}) operacionaliza precisamente esta distinción.

\subsection{Mecanismos de pago como señales de output}

Los mecanismos de financiación descritos en la Sección~\ref{sec:financiacion} tienen una lectura directa en términos de la teoría de contratos. El presupuesto global retrospectivo equivale, en la terminología de los incentivos, a un \textit{pago fijo} desvinculado del output: el hospital recibe recursos independientemente de su actividad o de la calidad de su desempeño. El pago por actividad ---sea por GRD, por módulos de concierto o por capitación ajustada--- equivale a un \textit{pago variable} vinculado a una señal observable del output.

El problema analítico central es que las señales disponibles no son igualmente precisas para todas las dimensiones del output. El volumen de actividad (altas, estancias, intervenciones) es relativamente fácil de medir y verificar. La intensidad de la atención ---la adecuación de los procedimientos diagnósticos y terapéuticos a la situación clínica del paciente--- es mucho más costosa de observar, más difícil de codificar en indicadores contractuales y más susceptible de manipulación. En la notación del modelo teórico, esta asimetría se captura mediante las varianzas de las señales de output: la señal de cantidad ($\sigma_1^2$) es relativamente precisa, mientras que la señal de intensidad ($\sigma_2^2$) es ruidosa.

La consecuencia, bien establecida en la teoría de multitarea \citep{holmstrom1991}, es que vincular los incentivos a la señal más precisa (cantidad) puede generar un desplazamiento de esfuerzo desde la dimensión menos verificable (intensidad). Este mecanismo es el fundamento teórico de la predicción central de esta tesis: la descentralización, al ir asociada a mecanismos de pago más vinculados a la actividad, puede mejorar la eficiencia en volumen pero deteriorar ---o al menos no mejorar proporcionalmente--- la eficiencia en intensidad.

\subsection{El parámetro $\psi_{12}$ y la naturaleza del servicio}

La última pieza de la traducción concierne al grado de interacción entre las tareas del médico. El Capítulo~\ref{ch:juridico} ha descrito que los hospitales prestan servicios de naturaleza heterogénea: servicios quirúrgicos, donde el acto de intervenir genera per se actividad diagnóstica y terapéutica asociada (preoperatorio, postoperatorio, seguimiento), y servicios médicos no quirúrgicos, donde la atención a un nuevo paciente compite en tiempo con la profundidad diagnóstica dedicada a cada uno.

En el modelo teórico, esta distinción se formaliza mediante el parámetro $\psi_{12}$, que mide la complementariedad o sustituibilidad entre las dos tareas del médico (producir cantidad y producir intensidad). En los servicios quirúrgicos, las tareas tienden a ser complementarias ($\psi_{12} < 0$): operar más implica per se más intensidad. En los servicios médicos, las tareas tienden a ser sustitutas ($\psi_{12} > 0$): ver más pacientes compite con el tiempo diagnóstico por paciente.

Esta correspondencia entre la naturaleza del servicio clínico y el parámetro de interacción entre tareas es la que justifica dos de los diseños empíricos de la tesis: el Diseño B, que estima fronteras separadas para servicios quirúrgicos y médicos, y el Diseño C, que explota la variación entre tipos de servicio dentro del mismo hospital.

\subsection{Síntesis de correspondencias}

La tabla siguiente resume las correspondencias entre las categorías jurídico-institucionales del Capítulo~\ref{ch:juridico} y los parámetros del modelo de agencia que se formaliza en el Capítulo~\ref{ch:modelo}:

\begin{table}[htbp]
\centering
\caption{Correspondencias entre categorías institucionales y parámetros del modelo de agencia}
\label{tab:correspondencias}
\begin{tabular}{p{5.5cm}p{5.5cm}}
\toprule
\textbf{Categoría jurídico-institucional} & \textbf{Parámetro del modelo económico} \\
\midrule
Forma de gestión (directa / indirecta) & Estructura de agencia (centralizada / descentralizada) \\[3pt]
Régimen del personal (estatutario / laboral) & Capacidad del hospital para diseñar el contrato del médico \\[3pt]
Mecanismo de financiación (presupuesto / actividad / capitación) & Tipo de señal de output y potencia de los incentivos \\[3pt]
Precisión de los indicadores de control & Varianzas de las señales ($\sigma_1^2$, $\sigma_2^2$) \\[3pt]
Naturaleza del servicio (quirúrgico / médico) & Parámetro de interacción entre tareas ($\psi_{12}$) \\[3pt]
Autonomía del gestor & Derechos residuales de decisión \\
\bottomrule
\end{tabular}
\end{table}

Estas correspondencias no pretenden agotar la complejidad institucional descrita en las secciones precedentes, sino identificar los ejes analíticos a lo largo de los cuales la diversidad organizativa del sistema hospitalario español genera variación relevante para un análisis de incentivos y eficiencia. El capítulo siguiente formaliza esta estructura en un modelo de agencia multitarea con tres agentes y dos outputs, del cual se derivan predicciones contrastables sobre cómo la estructura organizativa y el mecanismo de pago afectan de forma asimétrica a la eficiencia en cada dimensión del output hospitalario.


% ═══════════════════════════════════════════════════════════
% SECCIÓN 2.7 — TEXTO LITERAL (NO MODIFICAR)
% ═══════════════════════════════════════════════════════════

\section{Conclusión: pluralismo organizativo, información asimétrica y necesidad de evaluación empírica}

El sistema sanitario español admite, por habilitación legal explícita, una pluralidad de formas organizativas para la gestión hospitalaria. Esa pluralidad no es meramente formal: altera de modo sustantivo la autonomía del gestor, el régimen de contratación y retribución del personal, la distribución de riesgos entre financiador y proveedor, y los instrumentos disponibles para supervisar y disciplinar el desempeño.

La descripción institucional de este capítulo permite identificar tres rasgos del entorno hospitalario español que justifican un enfoque analítico de agencia multitarea bajo información asimétrica. Primero, la cadena de provisión es multiagente: financiador, hospital y médico operan en niveles distintos con información y objetivos parcialmente divergentes. Segundo, el output hospitalario es multidimensional: la legislación define las prestaciones en términos que incluyen cantidad, calidad, adecuación y continuidad, y no todas estas dimensiones son igualmente observables ni verificables. Tercero, la estructura organizativa ---centralizada o descentralizada--- determina quién diseña el contrato del médico y con qué instrumentos, lo que tiene consecuencias directas sobre la asignación del esfuerzo entre tareas.

Estos tres rasgos configuran un problema que ni el análisis puramente jurídico ni el análisis económico estándar de eficiencia pueden resolver por separado. El análisis jurídico describe las reglas del juego pero no predice sus consecuencias sobre el desempeño. El análisis de eficiencia convencional ---con un único indicador--- agrega dimensiones que pueden moverse en direcciones opuestas y pierde información esencial. La contribución de esta tesis es articular ambas perspectivas: el modelo teórico del Capítulo~\ref{ch:modelo} formaliza la estructura institucional como un problema de incentivos con soluciones analíticas y predicciones contrastables; la estimación empírica de los Capítulos~\ref{ch:resultados} y~\ref{ch:robustez} evalúa esas predicciones con datos del sistema hospitalario español; y las implicaciones de política del Capítulo~\ref{ch:politica} devuelven los resultados al terreno del diseño institucional.
